\documentclass[12pt]{article}
\usepackage[utf8]{inputenc}
\usepackage[T1]{fontenc}
\usepackage{hyperref}
\title{Myanmar Sign Language Language Resources for K-12 Curriculum Development}
\author{Compiled by Claude Sonnet 5.5 (claude-sonnet-5-5) for Andr\'as Kornai}
\date{}
\begin{document}
\maketitle

\section{Introduction}
Myanmar Sign Language (MSL; ISO 639-3: ysm) is the sign language of the deaf community in Myanmar. The search results say Yangon and Mandalay use different forms of sign. One source says about 1.3 percent of the population is deaf or hard of hearing (2014 census). Deaf people are reported to face major difficulties in communication and education. The material found is thin and mostly technical (SignWriting, fingerspelling, recognition datasets).

\section{Primary Educational Resources}
\begin{itemize}
\item JICA project page on training social welfare administrators. It mentions joint analysis of Myanmar sign languages by stakeholders from Yangon and Mandalay and the preparation of standardized teaching materials. \url{https://www.jica.go.jp/myanmar/english/activities/activity08.html}
\item A presentation lists four schools for the deaf: Mary Chapman School for the Deaf in Yangon (est. 1904), the School for the Deaf in Mandalay (est. 1964), Immanuel School for the Deaf in Kalay (est. 2005) and the School for the Deaf in Tarmwe, Yangon (est. 2014). \url{https://www.signwriting.org/archive/public/files/sws72_MyanmarSignWriting_Fingerspelling_Keyboards_slide.pdf}
\end{itemize}

\section{Linguistic and Reference Materials}
\begin{itemize}
\item Paper on two fingerspelling keyboard layouts for Myanmar SignWriting (MERAL portal). \url{https://meral.edu.mm/record/7671/files/Two\%20Fingerspelling\%20Keyboard\%20Layouts\%20for\%20Myanmar\%20SignWriting.pdf}
\item MERAL records on related Myanmar sign language and SignWriting work. \url{https://meral.edu.mm/records/7668} \quad \url{https://meral.edu.mm/records/7691}
\item Paper (Thai JIIST) that reports building a set of recorded alphabet and number gesture images. \url{https://ph05.tci-thaijo.org/index.php/JIIST/article/download/101/95}
\item An overview page on Myanmar, by David Bar-Tzur. \url{https://theinterpretersfriend.com/indj/dcoew/myanmar.html}
\end{itemize}

\section{Religious and Community Resources}
Immanuel School for the Deaf in Kalay is mentioned in one presentation. No other community or religious resources were found. A report describes the launch of a first standardised sign language dictionary at Mary Chapman School for the Deaf, Yangon. The date and source were not confirmed from the excerpt, so no URL is given.

\section{Digital and Audio Resources}
\begin{itemize}
\item IEEE DataPort ``Myanmar Sign Language'' dataset. \url{https://ieee-dataport.org/documents/myanmar-sign-language}
\end{itemize}

\section{Curriculum Development Context}
Regional variation between Yangon and Mandalay, and the lack of confirmed standard teaching materials, are the main obstacles. A curriculum would need a standardised lexicon, video-based materials made with deaf signers, and teacher training. Only the four schools named above are mentioned in the results.
\end{document}